\scanpage{81}
\diagram{\begin{tikzcd}[column sep=small]
\mathrm H^{\unsure{*|2}}(S,\operatorname{End}(M)) \arrow[r,dash] & \mathrm H^{\unsure{*|2}}(S,\operatorname{End}(M)\otimes_\Lambda\underline{\mathrm O}_S)\simeq\mathrm H^2(S,\operatorname{End}(E))
\end{tikzcd}}
Enfin, sous les conditions de 1.5, $M_0$ étant aussi plat sur $\Lambda$, on aura que $\unsure{\xi|\zeta}_\Lambda$ est l'image d'un élément
\[
\unsure{\xi|\zeta}_0\in\mathrm H^2(B_\pi,\operatorname{End}(M_0)).
\]
(NB Pour calculer $\unsure{\xi|\zeta}_0$, par fonctorialité en $\pi$, on est ramené aussitôt au cas ``universel'' $\pi=\operatorname{Aut}(M_0)$, et on peut évidemment supposer $M_0=\Lambda^d=(\mathbb Z/n\mathbb Z)^d$, donc $\pi=\operatorname{GL}(d,\mathbb Z/n\mathbb Z)$ ....) D'après \uline{Illusie}, c'est l'obstruction à remonter la représentation de $\pi$ dans $M_0$ en une représentation dans un module libre sur $\mathbb Z/\unsure{p^2|n^2}\mathbb Z$.

\inserted{Lorsque $n=p$ (premier),} Elle est évidemment nulle si $M_0$ est de rang $\leq1$, mais sans doute pas nulle si $M_0$ de rang $d\geq2$, dans le cas universel $\pi=\operatorname{GL}(d,\mathbf F_p)$.

Notons que si le rev. $S'$ de $S$ est le rev. universel, alors
\diagram{\begin{tikzcd}
\mathrm H^2(B_\pi,-) \arrow[r] & \mathrm H^2(S,-)
\end{tikzcd}}
est \uline{injectif}, donc on trouvera un élément non nul \deleted{dans} $\unsure{\xi|\zeta}_\Lambda\in\mathrm H^2(S,\operatorname{End}(M))$.
\mspage{80}
